% Part: second-order-logic
% Chapter: sol-and-set-theory
% Section: cardinalities

\documentclass[../../../include/open-logic-section]{subfiles}

\begin{document}

\olfileid{sol}{set}{crd}

\olsection{أعداد المجموعات الأصلية}

\begin{explain}
كما عبّرنا عن تناهي المجال ولا نهائيته، وعن كونه !!{enumerable}
أو~!!{nonenumerable}، فلنا أن نعرّف للمجموعة الجزئية
من~$\Domain{{\OLId{latin-looped}{M}}}$ هذه الخصائص نفسها: التناهي واللا نهائية،
وكونها !!{enumerable} أو~!!{nonenumerable}.
\end{explain}

\begin{prop}
تكون الصيغة~$\fn{Inf}({\OLId{latin-looped}{X}}) \ident$
\begin{multline*}
\lexists[{\OLId{latin-ordinary}{u}}][(\lforall[{\OLId{latin-ordinary}{x}}][({\OLId{latin-looped}{X}}({\OLId{latin-ordinary}{x}}) \lif {\OLId{latin-looped}{X}}({\OLId{latin-ordinary}{u}}({\OLId{latin-ordinary}{x}})))] \land {}\\
  \lforall[{\OLId{latin-ordinary}{x}}][\lforall[{\OLId{latin-ordinary}{y}}][(\eq[{\OLId{latin-ordinary}{u}}({\OLId{latin-ordinary}{x}})][{\OLId{latin-ordinary}{u}}({\OLId{latin-ordinary}{y}})] \lif 
      \eq[{\OLId{latin-ordinary}{x}}][{\OLId{latin-ordinary}{y}}])]] \land {} \\\lexists[{\OLId{latin-ordinary}{y}}][({\OLId{latin-looped}{X}}({\OLId{latin-ordinary}{y}}) \land \lforall[{\OLId{latin-ordinary}{x}}][({\OLId{latin-looped}{X}}({\OLId{latin-ordinary}{x}})
      \lif \eq/[{\OLId{latin-ordinary}{y}}][{\OLId{latin-ordinary}{u}}({\OLId{latin-ordinary}{x}})])])])]
\end{multline*}
مرضاة مع تعيين المتغيرات~${\OLId{latin-ordinary}{s}}$ إذا وفقط إذا كانت~${\OLId{latin-ordinary}{s}}({\OLId{latin-looped}{X}})$ لانهائية.
\end{prop}

\begin{prop}
تكون الصيغة~$\fn{Count}({\OLId{latin-looped}{X}}) \ident $
\begin{multline*}
(\lnot\lexists[{\OLId{latin-ordinary}{x}}][{\OLId{latin-looped}{X}}({\OLId{latin-ordinary}{x}})] \lor {}\\
\lexists[{\OLId{latin-ordinary}{z}}][\lexists[{\OLId{latin-ordinary}{u}}][({\OLId{latin-looped}{X}}({\OLId{latin-ordinary}{z}}) \land 
    \lforall[{\OLId{latin-ordinary}{x}}][({\OLId{latin-looped}{X}}({\OLId{latin-ordinary}{x}}) \lif {\OLId{latin-looped}{X}}({\OLId{latin-ordinary}{u}}({\OLId{latin-ordinary}{x}})))] \land {} \\ \lforall[{\OLId{latin-looped}{Y}}][(({\OLId{latin-looped}{Y}}({\OLId{latin-ordinary}{z}}) \land
      \lforall[{\OLId{latin-ordinary}{x}}][({\OLId{latin-looped}{Y}}({\OLId{latin-ordinary}{x}}) \lif {\OLId{latin-looped}{Y}}({\OLId{latin-ordinary}{u}}({\OLId{latin-ordinary}{x}})))]) \lif {\OLId{latin-looped}{X}} \subseteq {\OLId{latin-looped}{Y}})])]])
\end{multline*}
مرضاة مع تعيين المتغيرات~${\OLId{latin-ordinary}{s}}$ إذا وفقط إذا كانت~${\OLId{latin-ordinary}{s}}({\OLId{latin-looped}{X}})$ !!{enumerable}.
\end{prop}

أثبتت مبرهنة كانتور وجود مجموعات !!{nonenumerable}، بل دلّت
على أن للأحجام اللانهائية مراتب مختلفة لا نهاية لعددها.
وتقيم نظرية المجموعات لهذه الأحجام حسابًا تامًا، فتجعل للمجموعات
أعدادًا أصلية لانهائية. أما المجموعات المنتهية فتقاس أحجامها
بالأعداد الطبيعية. وأول الأصليات اللانهائية أصلية المجموعات
!!{denumerable}، واسمها~$\aleph_{\OLMathNumeral{0}}$، أي «ألف-صفر».
وتليها~$\aleph_{\OLMathNumeral{1}}$، وهي أصغر حجم لا تكون المجموعة عنده قابلة
للتعداد، أي لا يكون حجمها~$\aleph_{\OLMathNumeral{0}}$.
ولتعريف كون~${\OLId{latin-looped}{X}}$ ذات حجم~$\aleph_{\OLMathNumeral{0}}$ نضع
$\fn{Aleph}_{\OLMathNumeral{0}}({\OLId{latin-looped}{X}}) \liff \fn{Inf}({\OLId{latin-looped}{X}}) \land \fn{Count}({\OLId{latin-looped}{X}})$.
وتكون~${\OLId{latin-looped}{X}}$ ذات حجم~$\aleph_{\OLMathNumeral{1}}$ متى كانت لانهائية، وكان كل جزء
منها إما منتهيًا، وإما من الحجم~$\aleph_{\OLMathNumeral{0}}$، وإما مكافئًا لها
عدديًا، ولم تكن هي من الحجم~$\aleph_{\OLMathNumeral{0}}$.
فعبارة ذلك الصورية هي
$\fn{Aleph_1}({\OLId{latin-looped}{X}}) \ident \fn{Inf}({\OLId{latin-looped}{X}}) \land
\lforall[{\OLId{latin-looped}{Y}}][({\OLId{latin-looped}{Y}} \subseteq {\OLId{latin-looped}{X}} \lif
  (\lnot\fn{Inf}({\OLId{latin-looped}{Y}}) \lor \fn{Aleph}_{\OLMathNumeral{0}}({\OLId{latin-looped}{Y}}) \lor \cardeq{{\OLId{latin-looped}{X}}}{{\OLId{latin-looped}{Y}}}))] \land \lnot
\fn{Aleph}_{\OLMathNumeral{0}}({\OLId{latin-looped}{X}})$. وعلى هذا النحو يعرّف الحجم~$\aleph_{\OLMathNumeral{2}}$،
ثم ما يليه.

ومن الأحجام ما يخصه النظر، وهو أصلية المتصل:
حجم~$\Pow{\Nat}$، أو حجم~$\Real$، فهما سواء.
ووصف مجموعة بأنها من حجم المتصل ممكن في منطق الرتبة الثانية
أيضًا، غير أنه يحتاج إلى مزيد بيان.

\end{document}
